\chapter{Environment Setup Reference}
\index{environment setup}\index{Atlas setup}\index{installation}%
This appendix consolidates every setup decision the book makes into one reference: the Atlas
hierarchy and cluster tiers, the local Community Server footprint, the client tools each chapter
assumes, and the cost guardrails that keep a training environment free. Chapter~0 walks through
each item with screenshots-level detail; use this appendix when you need the fact, not the tour.

\section{The Atlas Hierarchy: Organization, Project, Cluster}
\index{Atlas organization}\index{Atlas project}%
Atlas nests resources in three levels, and every access-control and billing decision attaches to
one of them.

\begin{center}
\footnotesize
\begin{tabular}{@{}p{2.6cm}p{4.6cm}p{5.6cm}@{}}
\toprule
\textbf{Level} & \textbf{Owns} & \textbf{Book convention} \\
\midrule
Organization & Billing, API keys, SSO/SCIM, all projects & One organization per learner or team \\
Project & Clusters, database users, network access, alerts, integrations & One project: \texttt{MongoDB Book Labs} \\
Cluster & Data, backups, tiers, regions & One M0 named \texttt{Cluster0} for cloud labs \\
\bottomrule
\end{tabular}
\end{center}

\noindent Confirm the intended project is selected in the console before creating anything: most
``where did my cluster go'' incidents are ``created in the wrong project'' incidents.

\section{Cluster Tiers and When the Book Uses Each}
\index{cluster tier}\index{M0}\index{Atlas tiers}%
\begin{center}
\footnotesize
\begin{tabular}{@{}p{2.2cm}p{3.4cm}p{6.2cm}@{}}
\toprule
\textbf{Tier} & \textbf{Cost posture} & \textbf{Characteristics and book usage} \\
\midrule
M0 (Free) & Free forever, no card & 512 MB storage, shared RAM, shared infrastructure. Chapters 0, 2--3, 9--11, 13--14 functional labs. Not for benchmarking. \\
Flex & Low cost, capped & Up to 5 GB; bridging tier between M0 and dedicated. Optional for this book. \\
M10+ (Dedicated) & Billed hourly & Dedicated RAM/CPU, storage autoscaling, backups with PITR, private endpoints, x.509, auditing, Online Archive. Required for Chapters 17--20's paid-feature labs. \\
Serverless & Billed per read/write unit & Minimal configuration; auto-scaling. Alternative for Chapter 16's serverless exercises. \\
Local Community & Free & Full engine access on Windows; required for Chapters 1, 5--8, 12 (storage files, replica sets, sharding, profiler). \\
\bottomrule
\end{tabular}
\end{center}

\begin{pitfall}
\textbf{Accidentally choosing a paid tier.} The cluster-creation screen places tiers near each
other. Verify the tier reads \texttt{M0}, the price shows free, and no credit-card step appears.
If the page asks about disk scaling, backup policy, or dedicated hardware, re-check the tier.
\end{pitfall}

\section{Local Community Server on Windows}
\index{MongoDB Community Server}\index{mongod.cfg}\index{Windows service}%
The MSI installer creates a Windows service named \texttt{MongoDB}, a YAML configuration file
\texttt{mongod.cfg}, a data directory (\texttt{storage.dbPath}), and a log file
(\texttt{systemLog.path}) under \texttt{C:\textbackslash Program Files\textbackslash
MongoDB\textbackslash Server\textbackslash <version>\textbackslash}. The settings that matter for
every local lab:

\begin{itemize}
    \item \texttt{net.bindIp: 127.0.0.1} --- localhost only; the correct default for all book labs.
    \item \texttt{net.port: 27017} --- the default; replica-set labs use 28017--28019 (Chapter 5) and higher.
    \item \texttt{storage.dbPath} --- where WiredTiger files live; inspect it in Chapter 1.
    \item \texttt{systemLog.path} --- the first place to look when the service fails to start.
\end{itemize}

\noindent Manage the service from PowerShell: \texttt{Get-Service -Name MongoDB} to check,
\texttt{Start-Service}/\texttt{Stop-Service -Name MongoDB} (elevated) to control. YAML indentation
uses spaces, never tabs; a rejected config writes its reason to the log path.

\section{Client and Data Tools}
\index{mongosh}\index{MongoDB Compass}\index{Database Tools}\index{PyMongo}%
\begin{center}
\footnotesize
\begin{tabular}{@{}p{3.2cm}p{4.6cm}p{4.8cm}@{}}
\toprule
\textbf{Tool} & \textbf{Role in this book} & \textbf{Verify with} \\
\midrule
\texttt{mongosh} & Primary CLI for every chapter & \texttt{mongosh --version} \\
MongoDB Compass & Visual inspection, schema sampling, index view & Connect with a saved lab URI \\
Database Tools & \texttt{mongodump}, \texttt{mongorestore}, \texttt{mongoimport}, \texttt{mongoexport} & \texttt{mongodump --version} \\
PyMongo & Python driver for migration, validation, and QE labs & \texttt{python -m pip show pymongo} \\
Atlas CLI & Scriptable Atlas administration (Chapters 16, 20--21) & \texttt{atlas --version} \\
Terraform + \texttt{mongodbatlas} provider & Declarative Atlas infrastructure (Chapter 21) & \texttt{terraform version} \\
\texttt{kubectl} + Minikube & Operator labs (Chapter 22) & \texttt{minikube version} \\
\bottomrule
\end{tabular}
\end{center}

\noindent The Database Tools install separately from \texttt{mongosh}: installing the shell does
not provide dump/restore/import/export. Add each tool's \texttt{bin} directory to \texttt{PATH}
and verify before starting the chapter that needs it.

\section{Drivers and Connection Hygiene}
\index{driver}\index{connection string}%
The Atlas \textbf{Connect} dialog emits per-target URIs (shell, Compass, driver). The SRV form,
\texttt{mongodb+srv://user:<password>@cluster0.example.mongodb.net/db?retryWrites=true\&w=majority},
discovers hosts through DNS. Rules the book enforces from Chapter 0 onward:

\begin{itemize}
    \item Never commit a URI containing credentials; inject secrets from environment variables or a CI secret store (Chapter 21).
    \item URL-encode passwords containing \texttt{@}, \texttt{:}, \texttt{/}, \texttt{?}, or \texttt{\#}.
    \item In PyMongo, construct one long-lived \texttt{MongoClient} per process with \texttt{serverSelectionTimeoutMS} set; ping before first use so configuration errors fail fast.
    \item Serverless functions must warm-scope the client and bound \texttt{maxPoolSize}, or cold starts exhaust cluster connections (Chapter 16).
\end{itemize}

\section{Terraform Provider Quick Start}
\index{Terraform}\index{mongodbatlas provider}%
Chapter 21 manages Atlas declaratively. The minimal configuration: provider
\texttt{mongodb/mongodbatlas} authenticated with an Atlas \emph{programmatic API key} (public and
private key pair created at the organization or project level), resources
\texttt{mongodbatlas\_project}, \texttt{mongodbatlas\_cluster},
\texttt{mongodbatlas\_database\_user}, and \texttt{mongodbatlas\_project\_ip\_access\_list}. Keep
the API private key in the CI secret store---never in a committed \texttt{.tfvars}---and run
\texttt{terraform plan} as a pull-request gate before any \texttt{apply}
\cite{mongodbTerraformDocs,terraformDocs}.

\section{Cost Guardrails}
\index{cost guardrails}\index{billing alerts}%
\begin{itemize}
    \item Create exactly one M0 cluster for the book; load only small synthetic datasets.
    \item When M0 limits bite (512 MB storage, shared RAM), drop test data or move to the local server before considering an upgrade.
    \item If any lab requires a paid tier (Chapters 17--20), document expected cost, duration, cleanup step, and owner \emph{before} enabling it; pause or terminate the cluster when the lab ends.
    \item Configure Atlas billing alerts the moment a payment method exists; in an enterprise, wire the project into the organization's cost-management process.
    \item Local labs have their own quota risks: a large import can fill the workstation disk, and noisy workloads grow log files. Use small datasets and clean up per chapter.
\end{itemize}

\section{Per-Chapter Tool Requirements}
\index{tool requirements}%
\begin{center}
\footnotesize
\begin{tabular}{@{}p{1.2cm}p{5.2cm}p{6.0cm}@{}}
\toprule
\textbf{Ch.} & \textbf{Environment} & \textbf{Additional tools} \\
\midrule
0 & Atlas M0 + local Community Server & mongosh, Compass, Database Tools, PyMongo \\
1 & Local server (WiredTiger files, config) & mongosh, PowerShell \\
2--3 & Atlas M0 or local & mongosh, PyMongo (migration/validation scripts) \\
4 & Local PostgreSQL + Atlas M0 & Relational Migrator (or mapping file), PyMongo \\
5 & Three local \texttt{mongod} instances & mongosh, Windows firewall for partition drills \\
6--7 & Local sharded cluster (mongos + CSRS + shards) & mongosh \\
8 & Local sharded cluster or Atlas multi-region & mongosh \\
9 & Local server or Atlas & mongosh, \texttt{explain()}, data generator \\
10--11 & Atlas M0 or local & mongosh \\
12 & Local server & mongosh, profiler, \texttt{serverStatus} \\
13 & Atlas (Search supported from M0) & Atlas console or CLI for index definitions \\
14 & Atlas (Vector Search from M0) & Embedding provider account, PyMongo \\
15 & Atlas Stream Processing instance & Kafka (or Atlas-as-source), mongosh \\
16 & Atlas & Atlas CLI/Triggers, Data API, function runtime \\
17 & Local Enterprise or Atlas M10+ & OpenSSL for x.509 certificates \\
18 & Atlas M10+ or Enterprise 7.0+ & PyMongo + QE shared library, AWS KMS (or local key provider for learning) \\
19 & Atlas M10+ & Private endpoint-capable cloud account, SIEM or file sink \\
20 & Atlas M10+ & Atlas CLI, Data Federation \\
21 & Atlas + Git repository & Terraform, \texttt{mongodbatlas} provider, CI \\
22 & Minikube on Windows & kubectl, MongoDB Community Operator manifests \\
23 & Atlas (backups) or Ops Manager & mongodump/mongorestore, runbook template \\
24 & All of the above & Timed-evaluation protocol from the chapter \\
\bottomrule
\end{tabular}
\end{center}

\begin{tipbox}
Re-run the Chapter 0 verification sequence---\texttt{db.runCommand(\{ ping: 1 \})}, insert, find,
and \texttt{db.serverStatus().storageEngine}---after any OS update, Atlas maintenance event, or
tool upgrade. A two-minute transcript beats a broken Thursday lab.
\end{tipbox}
